\documentclass[11pt]{article}
\usepackage[margin=1in]{geometry}
\usepackage{hyperref}
\usepackage{xcolor}
\usepackage{longtable}
\usepackage{array}
\usepackage{enumitem}
\definecolor{goalosblue}{HTML}{0B1020}
\definecolor{goalosgold}{HTML}{B78A00}
\hypersetup{colorlinks=true, linkcolor=goalosblue, urlcolor=goalosblue}
\title{\textbf{AEP-003 --- ProofPacket Schema}\\\large Atomic Proof Units for Machine Work}
\author{Vincent Boucher\\QUEBEC.AI \& MONTREAL.AI}
\date{Version 1.1 Institutional Edition --- 2026-06-05}
\begin{document}
\maketitle

\begin{abstract}
AEP-003 defines the ProofPacket: the atomic evidence unit emitted by an AI-agent run, machine-work step, tool call, evaluation, policy decision, selection decision, rollout event, rollback event, or public-safe report event. AEP-001 defines the protocol. AEP-002 defines the Evidence Docket. AEP-003 defines the packet-level schema that makes machine work traceable, portable, tamper-evident, and composable.
\end{abstract}

\section*{Canonical Thesis}
A model can answer. An agent can act. An institution must prove. AEP-003 makes proof atomic.

\section*{Canonical Law}
No proof, no evolution.\\
No eval, no propagation.\\
No rollback, no release.

\section{Required Fields}
A conforming ProofPacket must include packet\_id, schema, schema\_version, packet\_type, created\_at, producer, docket\_id, commitment\_id, run\_id, claim\_refs, evidence\_refs, boundary, payload, hash, and claim\_boundary.

\section{Packet Types}
\begin{longtable}{p{0.25\linewidth}p{0.65\linewidth}}
\textbf{Packet Type} & \textbf{Purpose}\\
commit & Records mission, authorization, constraints, success criteria, and failure criteria.\\
trace\_event & Records a relevant execution event without necessarily publishing the full trace.\\
tool\_call & Records tool use, permission class, approval state, and result summary.\\
policy\_decision & Records policy decision, denial, escalation, or approval requirement.\\
approval & Records human or institutional authorization.\\
eval\_result & Records evaluation, test, check, score, or evaluator attestation.\\
evidence\_ref & Records a reference to evidence stored elsewhere.\\
risk\_event & Records a risk, incident, mitigation, or residual concern.\\
cost\_event & Records cost, latency, token count, compute time, or human-review burden.\\
selection\_decision & Records a Selection Gate decision.\\
rollout\_event & Records canary scope, monitoring plan, and propagation limit.\\
rollback\_event & Records rollback trigger, target, action, owner, and recovery.\\
public\_report & Records a public-safe report publication event.\\
\end{longtable}

\section{Canonical JSON and Hashing}
A ProofPacket should be hashed using deterministic canonical JSON: remove mutable signature fields, set hash to a zero-placeholder, sort object keys, encode as UTF-8, and hash using SHA-256.

\[
packet\_hash = SHA256(canonical\_json(packet\_without\_signature))
\]

\section{Conformance Levels}
\begin{longtable}{p{0.20\linewidth}p{0.70\linewidth}}
Level 0 & Informal human-readable proof fragment.\\
Level 1 & Valid JSON packet.\\
Level 2 & Hashable packet with valid canonical hash.\\
Level 3 & Docket-linked packet.\\
Level 4 & Institutional packet with policy, eval, risk, cost, boundary, attestation, and Selection Gate references.\\
Level 5 & Sovereign / regulated packet with jurisdiction, retention, protected-boundary classification, signer, timestamp authority, and restricted appendix controls.\\
\end{longtable}

\section{Claim Boundary}
AEP-003 does not claim achieved AGI, achieved ASI, perfect safety, legal compliance certification, financial or legal advice, guaranteed ROI, production readiness, government endorsement, or national-security readiness. It defines an atomic proof schema.

\section*{Canonical Public Line}
AEP-001 defines the protocol. AEP-002 defines the docket. AEP-003 defines the packet. GoalOS makes machine work provable one packet at a time.

\end{document}
